\section{Computation Budget Guidance}
\label{sec:budget}

\subsection{Overview}

The 15-algorithm portfolio spans four orders of magnitude in runtime:
\flip{} on NH completes in 0.3 s; \sa{} on TX requires 82.1 s.
For a 50-state production sweep using the recommended configurations,
practitioners need a compute budget estimate and a parallelism strategy.
This section provides both.

\subsection{Per-State Runtime Estimates}

Table~\ref{tab:budget-per-state} gives estimated wall-clock runtime per state
for the four most common production configurations.
Estimates are derived from the five benchmark states (NC, WI, TX, FL, NH)
by linear interpolation in $k \cdot N$ (the dominant cost factor for
METIS-class algorithms).

\begin{table}[ht]
\centering
\caption{%
  Estimated wall-clock runtime per state (seconds) for four production
  configurations.
  Estimates for $k, N$ values not in the benchmark states are interpolated.
  ``16w'' indicates \texttt{--workers 16}; all others \texttt{--workers 1}.
  These are rough estimates ($\pm 30\%$) for planning purposes.
}
\label{tab:budget-per-state}
\small
\begin{tabular}{@{}lrrrrr@{}}
\toprule
Configuration & NC ($k=14$) & WI ($k=8$) & TX ($k=38$) & FL ($k=28$) & Typical 50-state \\
\midrule
METIS + Convergence (default)
  &  $\sim$12 & $\sim$5 & $\sim$38 & $\sim$28 & $\sim$500 s \\
SA + Short-Burst (EC optimal)
  &  $\sim$40 & $\sim$18 & $\sim$135 & $\sim$98 & $\sim$1{,}800 s \\
CVD-Geo + SMC-Percentile (PP + calibrated)
  &  $\sim$27 & $\sim$12 & $\sim$95 & $\sim$72 & $\sim$1{,}100 s \\
SA + PT-16w (EC + large-$k$ optimal)
  &  $\sim$37 & $\sim$16 & $\sim$124 & $\sim$90 & $\sim$1{,}650 s \\
\bottomrule
\end{tabular}
\end{table}

The default configuration (METIS + Convergence sweep) completes a 50-state
congressional sweep in approximately 500 s ($\approx 8$ minutes) on a
single thread, or $\approx 30$ s with 16-thread Rayon parallelism.
This is the recommended configuration for exploratory runs and parameter tuning.

\subsection{Practical Algorithm Selection Under Budget Constraints}

\paragraph{Tight budget ($< 5$ minutes, single thread).}
Use METIS + Convergence (default) or BFS + ConvergenceSweep.
ILP is applicable only to NH; exclude from production sweep or apply
selectively.
Do not use SA, CVD-Geo, PT, or SMC-Percentile.

\paragraph{Moderate budget ($< 30$ minutes, single thread or $< 5$ minutes
with 16 threads).}
Use BFS or CVD-Geo structure + SMC-Percentile or Merge-Split search.
This configuration covers PP-maximising and calibrated-distribution use
cases at approximately $2\times$ the default runtime.

\paragraph{Full budget ($< 2$ hours, 16 threads).}
SA structure + PT search for EC-optimal results.
CVD-Geo structure + SMC-Percentile search for PP + calibrated results.
Adaptive Multi-scale for large-$k$ states (TX, FL, CA, NY) for mixing
quality.

\subsection{Parallelism Strategy}

The three-layer compositor exposes two levels of parallelism:

\paragraph{State-level parallelism.}
Each state is independent.
For a 50-state sweep, distribute states across nodes with
\texttt{--workers 1} per state and $\leq 50$ parallel jobs.
This is the most efficient parallelism strategy for homogeneous clusters.

\paragraph{Algorithm-level parallelism (\pt{} only).}
\pt{} parallelises internally across replicas.
With 8 replicas and \texttt{--workers 8} or \texttt{--workers 16},
PT achieves near-linear speedup up to 8 cores.
For a node with 16 cores handling TX (the bottleneck state in any sweep),
run 2 states in parallel each with 8 workers:
\texttt{--workers 8 --replicas 8}.

\paragraph{Anti-pattern to avoid.}
Do not parallelise chain steps within a single state (e.g., using multiple
threads for one Merge-Split chain).
The Markov chain is inherently sequential; multi-threading a single chain
provides no benefit and introduces synchronisation overhead.
Only \pt{} benefits from within-state multi-threading (replica parallelism).

\subsection{Which Algorithms Are Practical for 50-State Sweeps}

Table~\ref{tab:practical} classifies all 15 algorithms by production
practicality for a 50-state congressional sweep.

\begin{table}[ht]
\centering
\caption{%
  Production practicality for a 50-state congressional sweep.
  ``Full sweep'' means all 50 states; ``Select'' means large-$k$ states only
  or NH only; ``Not practical'' means runtime or hardware requirements
  preclude routine use.
}
\label{tab:practical}
\small
\begin{tabular}{@{}lll@{}}
\toprule
Algorithm & Production Practicality & Notes \\
\midrule
METIS        & Full sweep & Default; fastest \\
BFS Growth   & Full sweep & Slightly faster than METIS \\
CVD-Graph    & Full sweep & $2\times$ METIS runtime \\
CVD-Geo      & Full sweep & Requires centroid data; $2\times$ METIS \\
SA           & Full sweep & $10\times$ METIS; budget $\sim$30 min/sweep \\
ILP          & NH only    & Applicable to $k=2$ states; $60\times$ METIS \\
\midrule
ConvergenceSweep  & Full sweep & Default search; fast \\
PercentileSweep   & Full sweep & $\sim$1.5$\times$ Convergence \\
BisectionEnsemble & Full sweep & $\sim$2$\times$ Convergence \\
Merge-Split       & Full sweep & Recommended general-purpose \\
Forest ReCom      & Full sweep & Similar to Merge-Split \\
Short-Burst       & Full sweep & EC optimal, single-core \\
ShortBurstForest  & Full sweep & EC + MH; $\sim$50\% slower than Short-Burst \\
SMC               & Full sweep & Calibrated; $\sim$5$\times$ Merge-Split \\
SMC-Percentile    & Full sweep & Calibrated point estimate; $\sim$SMC cost \\
VRA-Aware         & Select     & VRA jurisdictions only \\
Flip              & Select     & Sensitivity analysis at a fixed plan \\
Multi-scale       & Large-$k$  & TX, FL, CA, NY, PA, OH; not needed for $k<15$ \\
Adaptive MS       & Large-$k$  & Preferred over fixed Multi-scale for $k \geq 20$ \\
Parallel Tempering & Large-$k$ (multi-core) & 8--16 cores; EC optimal for large-$k$ \\
\bottomrule
\end{tabular}
\end{table}

\subsection{Recommended 50-State Production Stack}

For a complete 50-state congressional sweep that balances quality and runtime:

\begin{enumerate}
  \item \textbf{Default pass}: All 50 states with
        \texttt{--structure prime-factor --weights county --search convergence}.
        Approximately 8 minutes single-threaded.
        This produces the official plan for all states.

  \item \textbf{Large-$k$ pass}: TX, FL, CA, NY, PA, OH ($k \geq 15$) with
        \texttt{--structure sa-bisect --search adaptive-multiscale --workers 16}.
        Approximately 40 minutes.
        Improves EC and ACF for the six states that most benefit from
        large-$k$ methods.

  \item \textbf{ILP pass}: NH only with
        \texttt{--structure ilp --search single}.
        Approximately 13 seconds.
        Certifies optimal EC for the two-district state.

  \item \textbf{VRA pass}: Jurisdictions with active Section~2 claims with
        \texttt{--structure ratio-optimal --weights vra-aligned --search vra-recom}.
        Approximately 10--20 minutes depending on jurisdiction count.
\end{enumerate}

Total estimated runtime for the complete four-pass stack on a 16-core machine:
approximately 60--90 minutes.
